\documentclass[11pt]{article}

% ─── Packages ───────────────────────────────────────────────────────────────
\usepackage[margin=1in]{geometry}
\usepackage{amsmath,amssymb,amsthm}
\usepackage{graphicx}
\usepackage{booktabs}
\usepackage{hyperref}
\usepackage{natbib}
\usepackage{algorithm}
\usepackage{algpseudocode}
\usepackage{xcolor}
\usepackage{bm}

\newtheorem{theorem}{Theorem}
\newtheorem{proposition}{Proposition}

\hypersetup{
    colorlinks=true,
    linkcolor=blue,
    citecolor=blue,
    urlcolor=blue,
}

% ─── Macros ─────────────────────────────────────────────────────────────────
\newcommand{\vech}{\bm{h}}
\newcommand{\anchor}[1]{A_{#1}}
\newcommand{\barrier}{\tau}
\newcommand{\dive}{D}
\newcommand{\nhat}{\hat{\bm{n}}}
\newcommand{\R}{\mathbb{R}}

\title{Iterative Attractor Dynamics for Classification:\\A Dynamical Classification Head with Anchor-Guided Basin Settling}

\author{
    Chetan Patil \\
    \texttt{chetan12patil@gmail.com} \\
    \url{https://github.com/chetanxpatil/livnium}
}

\date{March 2026}

\begin{document}
\maketitle

% ═══════════════════════════════════════════════════════════════════════════
\begin{abstract}
We present Livnium, a classification head implemented as a discrete-time
dynamical system. Instead of projecting an encoded representation to logits
via a single linear layer, Livnium evolves the hidden state through
$L$ steps of geometry-aware updates under learned anchor forces before
readout. Three anchor vectors define label basins in a shared embedding
space; the state vector $\vech$ moves under cosine-based divergence forces
until it settles near one basin, and the hidden state is iteratively
refined under anchor-guided forces before final readout. We define a
Lyapunov function $V(\vech) = (\barrier - \cos(\vech, \anchor{y}))^2$
and show that, in the absence of the learned residual, the dynamics are
locally contracting (a first-order local descent argument, not a global
stability proof). The approach is evaluated on the Stanford Natural
Language Inference (SNLI) benchmark, achieving 76\% dev accuracy with a
frozen bag-of-words encoder. We report the geometric properties of the
force field, including a measured 135.2\textdegree{} mismatch between
the implemented Euclidean radial forces and the true cosine gradient,
and discuss implications for convergence. Code and pretrained weights
are publicly available.
\end{abstract}

% ═══════════════════════════════════════════════════════════════════════════
\section{Introduction}
\label{sec:intro}

The standard neural network classifier computes $y = \text{softmax}(Wh + b)$:
a single affine transformation from a learned representation $h$ to class
logits. This is a static map --- the label is determined by a single
projection of the representation.

We explore an alternative: what if the classification head were a small
dynamical system that \emph{evolves} the representation through several
steps before producing a label? Instead of computing the label from the
raw encoding $h_0$, we let the state evolve:
\[
    h_0 \;\to\; h_1 \;\to\; \cdots \;\to\; h_L \;\to\; \text{logits}
\]
where each update step applies learned residuals and geometric forces
from label-specific anchor vectors. The final state $h_L$ is the result
of a settling process, not a one-shot projection.

This idea has historical roots. Hopfield networks \citep{hopfield1982}
used energy minimisation for content-addressable memory. Modern
equilibrium models \citep{bai2019deep} compute representations as fixed
points of implicit layers. Our approach sits between the two: it uses
an explicit, finite-step dynamical rule with interpretable geometric
structure, rather than an implicit fixed-point iteration.

The contribution of this paper is:
\begin{enumerate}
    \item A concrete update rule for attractor-based classification
          with three label basins and cosine-ring equilibria (Section~\ref{sec:method}).
    \item A Lyapunov analysis showing local contraction of the base
          dynamics (Section~\ref{sec:lyapunov}).
    \item An empirical characterisation of the geometric inconsistency
          between the implemented forces and true cosine gradients
          (Section~\ref{sec:geometry}).
    \item SNLI results with analysis of per-class behavior and
          inference speed (Section~\ref{sec:experiments}).
\end{enumerate}

We emphasize that this work replaces only the classification head.
The encoder (a bag-of-words mean pooling in our experiments) is
orthogonal to the attractor dynamics. The question we address is:
\emph{does iterative basin settling provide useful inductive bias
relative to a standard readout when the encoder is held fixed?}

% ═══════════════════════════════════════════════════════════════════════════
\section{Method}
\label{sec:method}

\subsection{Setup}

Given a premise $p$ and hypothesis $q$, an encoder produces sentence
vectors $\bm{v}_p, \bm{v}_q \in \R^d$. The initial state for the
collapse dynamics is:
\[
    h_0 = \bm{v}_q - \bm{v}_p + \epsilon, \quad \epsilon \sim \mathcal{N}(0, 0.01^2 I)
\]
The difference vector captures the directional relationship between
premise and hypothesis. The small noise $\epsilon$ breaks symmetry
in the attractor landscape.

\subsection{Label Geometry}

Three learned anchor vectors $\anchor{E}, \anchor{C}, \anchor{N} \in \R^d$
(entailment, contradiction, neutral) define the geometry of the label
space. Each anchor is unit-normalised during the forward pass. The
\emph{equilibrium} for each label is not the anchor point itself but a
\textbf{cosine-similarity ring}:
\[
    \cos(h, \anchor{y}) = \barrier
\]
where $\barrier = 0.38$ is a fixed barrier constant. This defines a
$(d-2)$-dimensional manifold on the unit hypersphere --- a ``ring''
at a fixed angular distance from each anchor.

\subsection{Update Rule}

At each step $t = 0, \ldots, L-1$:
\begin{equation}
\label{eq:update}
    h_{t+1} = h_t
    + \delta_\theta(h_t)
    - s_y \cdot \dive(h_t, \anchor{y}) \cdot \nhat(h_t, \anchor{y})
    - \beta \cdot B(h_t) \cdot \nhat(h_t, \anchor{N})
\end{equation}
where:
\begin{align}
    \dive(h, A) &= \barrier - \cos(h, A) \label{eq:divergence} \\
    \nhat(h, A) &= \frac{h - A}{\|h - A\|} \label{eq:direction} \\
    B(h) &= 1 - |\cos(h, \anchor{E}) - \cos(h, \anchor{C})| \label{eq:boundary}
\end{align}

The terms in \eqref{eq:update} are:
\begin{itemize}
    \item $\delta_\theta(h_t)$: a learned residual (2-layer MLP with Tanh).
    \item $s_y \cdot \dive \cdot \nhat$: the \textbf{anchor force}. When
          $\cos(h, A_y) < \barrier$, divergence is positive and the force
          pulls $h$ toward the anchor. When $\cos(h, A_y) > \barrier$,
          the force pushes $h$ away. The equilibrium ring acts as an
          attractor, not the anchor itself.
    \item $\beta \cdot B(h) \cdot \nhat(h, A_N)$: the \textbf{neutral
          boundary force}. $B(h)$ measures equidistance from $\anchor{E}$
          and $\anchor{C}$; when $h$ sits on the E--C boundary, $B \approx 1$
          and the force pulls $h$ toward neutral territory.
\end{itemize}

During training, only the anchor corresponding to the correct label $y$
exerts force ($s_y > 0$; other strengths are zero). During inference,
all three anchors act simultaneously with fixed hyperparameter strengths
($s_E = s_C = 0.1$, $s_N = 0.05$; these are the same values used during
training, not learned or adapted), and the label is determined by which
basin the state settles nearest.

A soft norm cap prevents unbounded growth:
\[
    h \leftarrow \begin{cases}
        h \cdot \frac{10}{\|h\| + 10^{-8}} & \text{if } \|h\| > 10 \\
        h & \text{otherwise}
    \end{cases}
\]

\subsection{Classification Head}

The full classification pipeline is: encoder $\to$ attractor dynamics
$\to$ readout network. After $L$ collapse steps, the refined state
$h_L$ is passed to a two-layer readout MLP along with auxiliary
geometric features: cosine alignment between $\bm{v}_p$ and $\bm{v}_q$,
opposition signal $\cos(-\bm{v}_p, \bm{v}_q)$, Euclidean distance
$\|\bm{v}_q - \bm{v}_p\|$, norm features, and a learned neutral
alignment signal. The readout MLP outputs logits for the three SNLI
classes. We note that the final label decision is made by this readout
network, not purely by basin membership; the attractor dynamics refine
the representation that the readout network receives.

\subsection{Dynamic Basin Field}

During training, a \textbf{basin field} maintains multiple micro-basins
per label. Each basin has a center vector updated via exponential moving
average on the unit sphere. New basins are spawned when tension (divergence
magnitude) stays high and alignment is low, indicating an underserved
region of the label manifold. Basins are periodically pruned and merged.

At inference, the basin field is read-only: the system routes each
sample to the nearest basin using label-free cosine similarity. No
ground-truth labels are used during inference routing.

% ═══════════════════════════════════════════════════════════════════════════
\section{Lyapunov Analysis}
\label{sec:lyapunov}

Define the candidate Lyapunov function:
\begin{equation}
\label{eq:lyapunov}
    V(h) = \big(\barrier - \cos(h, \anchor{y})\big)^2 = \dive(h, \anchor{y})^2
\end{equation}

$V = 0$ exactly on the equilibrium ring $\cos(h, A_y) = \barrier$.
$V$ serves as a local diagnostic measuring distance from the target
ring, not a global energy function for the full system.

\begin{proposition}[Local contraction without learned residual]
When $\delta_\theta = 0$ and only the correct anchor force acts,
$V(h_{t+1}) \leq V(h_t)$ for sufficiently small force strength $s_y$.
This is a first-order local descent argument, not a global stability
proof.
\end{proposition}

\begin{proof}[Proof sketch]
The anchor force component is $-s_y \cdot \dive(h, A) \cdot \nhat(h, A)$.
The change in cosine alignment due to this force involves:
\[
    \nabla_h \cos(h, A) \cdot \nhat(h, A)
    = -\frac{r \sin^2\theta}{\alpha \|h - A\|}
\]
where $\alpha = \|h\|$, $r = \|A\|$, and $\theta$ is the angle
between $h$ and $A$. Since $\sin^2\theta \geq 0$ and all norms are
positive, the dot product is non-positive. Therefore the anchor force
moves $\cos(h, A)$ toward $\barrier$ (reduces $|\dive|$), which
reduces $V$.
\end{proof}

Empirical verification (5,000 samples, $\delta_\theta$ scale varied):

\begin{table}[ht]
\centering
\begin{tabular}{@{}lcc@{}}
\toprule
$\delta_\theta$ scale & $V(h_{t+1}) \leq V(h_t)$ & mean $\Delta V$ \\
\midrule
0.001 & 100.0\% & $-0.00131$ \\
0.01  & 99.3\%  & $-0.00112$ \\
0.05  & 70.9\%  & $-0.00040$ \\
0.10  & 61.3\%  & $+0.00001$ \\
\bottomrule
\end{tabular}
\caption{Lyapunov descent rate as the learned residual strength increases.
The base dynamics are contracting; the MLP residual can partially oppose
contraction at larger scales.}
\label{tab:lyapunov}
\end{table}

\textbf{Limitation.} Global convergence is not proven. The learned
residual $\delta_\theta$ can dominate the anchor force at large scales,
and finite step sizes can overshoot the attractor ring. The Lyapunov
guarantee is local and conditional on $\delta_\theta$ being small
relative to the anchor force.

% ═══════════════════════════════════════════════════════════════════════════
\section{Geometric Inconsistency}
\label{sec:geometry}

A notable property of the update rule is that the force magnitude is
cosine-based (Eq.~\ref{eq:divergence}) but the force direction is
Euclidean radial (Eq.~\ref{eq:direction}). The true gradient of cosine
similarity with respect to $h$ is tangential on the hypersphere, not
radial.

To quantify this, we measured the angle between the implemented force
direction $\nhat(h, A)$ and the true cosine gradient direction
$\nabla_h \cos(h, A) / \|\nabla_h \cos(h, A)\|$ across 1,000 random
samples in $\R^{256}$:

\begin{equation}
    \text{mean angle} = 135.2^{\circ} \pm 2.5^{\circ}
\end{equation}

This means the system applies forces that are roughly orthogonal-plus-opposing
to the true gradient. In physical terms, this is a \textbf{non-conservative
force field} --- there is no scalar potential whose gradient produces
these forces.

Despite this, the system converges empirically. We interpret the anchor
force as a discrete-time attractor rule that happens to use cosine as a
distance proxy, not as gradient descent on an energy surface. The
Euclidean direction provides a consistent ``pull toward the anchor''
that, combined with the cosine-based magnitude, produces effective basin
capture even though the geometry is formally inconsistent.

This inconsistency may even be beneficial: the non-conservative
``sideways'' component could act as implicit regularisation, preventing
the state from collapsing directly onto the anchor and instead
encouraging exploration of the attractor ring.

% ═══════════════════════════════════════════════════════════════════════════
\section{Experiments}
\label{sec:experiments}

\subsection{Dataset and Encoder}

We evaluate on the Stanford Natural Language Inference (SNLI) dataset
\citep{bowman2015snli}, which contains approximately 570k sentence pairs
labeled as entailment, contradiction, or neutral.

The encoder is a pretrained bag-of-words model: a learned embedding table
(trained on WikiText-103 \citep{merity2016wikitext}) followed by mean
pooling over tokens. This produces 256-dimensional sentence vectors.
The encoder is frozen during attractor head training.

\subsection{Configuration}

\begin{table}[ht]
\centering
\begin{tabular}{@{}ll@{}}
\toprule
Parameter & Value \\
\midrule
Embedding dimension & 256 \\
Collapse steps $L$ & 6 \\
Barrier constant $\barrier$ & 0.38 \\
Anchor strengths $s_E, s_C$ & 0.1 \\
Anchor strength $s_N$ & 0.05 \\
Neutral boundary boost $\beta$ & 0.05 \\
Total parameters & ${\sim}2$M \\
\bottomrule
\end{tabular}
\caption{Model configuration.}
\label{tab:config}
\end{table}

\subsection{Results}

\begin{table}[ht]
\centering
\begin{tabular}{@{}lc@{}}
\toprule
Class & Dev Accuracy \\
\midrule
Overall & \textbf{76.00\%} \\
Entailment & 80.6\% \\
Contradiction & 75.2\% \\
Neutral & 72.2\% \\
\bottomrule
\end{tabular}
\caption{SNLI dev set accuracy (label-free inference routing).}
\label{tab:results}
\end{table}

For context, bag-of-words baselines typically achieve ${\sim}80\%$ on
SNLI, decomposable attention \citep{parikh2016decomposable} reaches
${\sim}86\%$, and fine-tuned BERT \citep{devlin2019bert} achieves
${\sim}90\%$.

Our result is below the bag-of-words baseline. However, the attractor
head receives mean-pooled bag-of-words vectors, which are known to
lose word order. Sentences such as ``a dog bites a man'' and ``a man
bites a dog'' produce nearly identical encodings, creating an
encoder-side ceiling that is difficult to separate from head-side
limitations. A controlled comparison with a frozen BERT encoder would
isolate the attractor head's contribution.

\subsection{Inference Speed}

\begin{table}[ht]
\centering
\begin{tabular}{@{}lcc@{}}
\toprule
Component & ms / batch (32) & Samples / sec \\
\midrule
Livnium attractor head & 0.4 & 85,335 \\
\bottomrule
\end{tabular}
\caption{Head-only inference speed on CPU. This measures only the
attractor dynamics and readout, excluding all encoding time. It is
\textbf{not} an end-to-end comparison: a fair speed comparison would
need to include the encoder for both systems.}
\label{tab:speed}
\end{table}

\subsection{Neutral Class Analysis}

Neutral is the weakest class (72.2\% vs.\ 80.6\% for entailment).
The confusion matrix shows that neutral samples are frequently
misclassified as both entailment and contradiction. The neutral boundary
force (Eq.~\ref{eq:boundary}) was specifically designed to address this,
pulling states toward neutral when they are equidistant from the E and C
anchors. Without it, neutral accuracy drops further.

The neutral class is intrinsically harder because it is defined by
the \emph{absence} of a strong entailment or contradiction signal.
With a bag-of-words encoder that loses word order, many genuinely
neutral pairs produce embeddings that overlap with entailment or
contradiction pairs.

% ═══════════════════════════════════════════════════════════════════════════
\section{Related Work}
\label{sec:related}

\textbf{Hopfield networks} \citep{hopfield1982} pioneered energy-based
attractor dynamics for memory retrieval. Modern Hopfield networks
\citep{ramsauer2020hopfield} extend this to continuous states with
exponential storage capacity.

\textbf{Deep equilibrium models} (DEQ) \citep{bai2019deep} compute
representations as fixed points of an implicit layer, using root-finding
rather than explicit unrolling. Livnium uses explicit $L$-step dynamics
with no fixed-point solver.

\textbf{Neural ODEs} \citep{chen2018neural} parameterise continuous-time
dynamics. Livnium operates in discrete time with interpretable geometric
forces rather than a learned vector field.

\textbf{Prototype networks} \citep{snell2017prototypical} classify by
distance to class prototypes. Livnium's anchors serve a similar role,
but the state is iteratively refined rather than directly compared.

\textbf{Natural language inference} baselines range from
bag-of-words \citep{bowman2015snli} to decomposable attention
\citep{parikh2016decomposable} to fine-tuned transformers
\citep{devlin2019bert}. Our work is orthogonal to the encoder and
addresses only the classification head.

% ═══════════════════════════════════════════════════════════════════════════
\section{Limitations and Future Work}
\label{sec:limitations}

\begin{enumerate}
    \item \textbf{Train/inference mismatch.} During training, only the
          correct-label anchor exerts force; during inference, all three
          anchors act simultaneously. This means the dynamics the model
          learns are not identical to the dynamics used at test time.
          The mismatch may limit settling accuracy and is a structural
          gap that future work should address.
    \item \textbf{Encoder ceiling.} The bag-of-words encoder cannot
          distinguish word order, which limits accuracy regardless of the
          classification head. A fair evaluation requires a controlled
          comparison: frozen BERT encoder with attractor head vs.\ frozen
          BERT encoder with linear head.
    \item \textbf{No global convergence.} The Lyapunov guarantee is local
          and conditional on the learned residual being small. Finite
          step sizes and the learned MLP can cause the state to escape
          basins.
    \item \textbf{Geometric inconsistency.} The 135\textdegree{} mismatch
          between Euclidean radial forces and true cosine gradients is
          acknowledged but not resolved. Whether replacing the direction
          with the true tangential gradient would improve or harm
          performance is an open question.
    \item \textbf{Limited benchmarking.} Results are reported on SNLI
          only. Evaluation on MultiNLI, adversarial NLI, and
          out-of-distribution robustness tasks would better characterise
          the approach.
    \item \textbf{Hyperparameter sensitivity.} The barrier constant
          $\barrier = 0.38$, force strengths, and number of collapse
          steps were tuned manually. A principled selection method
          is needed.
\end{enumerate}

% ═══════════════════════════════════════════════════════════════════════════
\section{Conclusion}

We have presented a classification head that operates as an iterative
dynamical system rather than a static projection. The core idea ---
replacing $h \to Wh + b$ with $h_0 \to h_1 \to \cdots \to h_L$ under
geometric anchor forces --- is simple and fast, and comes with a partial
stability guarantee via a Lyapunov function.

The current results do not exceed standard baselines on SNLI. However,
the approach opens a different design axis: classification heads with
built-in iterative refinement, interpretable trajectory dynamics, and
geometric stability properties. Whether this axis leads to practical
gains remains to be demonstrated with stronger encoders and broader
evaluation.

Code: \url{https://github.com/chetanxpatil/livnium} \\
Model: \url{https://huggingface.co/chetanxpatil/livnium-snli}

% ═══════════════════════════════════════════════════════════════════════════
\bibliographystyle{plainnat}
\begin{thebibliography}{10}

\bibitem[Bai et~al.(2019)]{bai2019deep}
S.~Bai, J.~Z. Kolter, and V.~Koltun.
\newblock Deep equilibrium models.
\newblock In \emph{NeurIPS}, 2019.

\bibitem[Bowman et~al.(2015)]{bowman2015snli}
S.~Bowman, G.~Angeli, C.~Potts, and C.~Manning.
\newblock A large annotated corpus for learning natural language inference.
\newblock In \emph{EMNLP}, 2015.

\bibitem[Chen et~al.(2018)]{chen2018neural}
R.~T.~Q. Chen, Y.~Rubanova, J.~Bettencourt, and D.~Duvenaud.
\newblock Neural ordinary differential equations.
\newblock In \emph{NeurIPS}, 2018.

\bibitem[Devlin et~al.(2019)]{devlin2019bert}
J.~Devlin, M.-W. Chang, K.~Lee, and K.~Toutanova.
\newblock {BERT}: Pre-training of deep bidirectional transformers for
  language understanding.
\newblock In \emph{NAACL-HLT}, 2019.

\bibitem[Hopfield(1982)]{hopfield1982}
J.~J. Hopfield.
\newblock Neural networks and physical systems with emergent collective
  computational abilities.
\newblock \emph{PNAS}, 79(8):2554--2558, 1982.

\bibitem[Merity et~al.(2016)]{merity2016wikitext}
S.~Merity, C.~Xiong, J.~Bradbury, and R.~Socher.
\newblock Pointer sentinel mixture models.
\newblock \emph{arXiv:1609.07843}, 2016.

\bibitem[Parikh et~al.(2016)]{parikh2016decomposable}
A.~Parikh, O.~T{\"a}ckstr{\"o}m, D.~Das, and J.~Uszkoreit.
\newblock A decomposable attention model for natural language inference.
\newblock In \emph{EMNLP}, 2016.

\bibitem[Ramsauer et~al.(2020)]{ramsauer2020hopfield}
H.~Ramsauer, B.~Sch{\"a}fl, J.~Lehner, P.~Seidl, M.~Widrich,
  T.~Adler, L.~Gruber, M.~Holzleitner, M.~Pavlovi{\'c}, G.~Sandve,
  V.~Greiff, D.~Kreil, M.~Kopp, G.~Klambauer, J.~Brandstetter,
  and S.~Hochreiter.
\newblock Hopfield networks is all you need.
\newblock In \emph{ICLR}, 2021.

\bibitem[Snell et~al.(2017)]{snell2017prototypical}
J.~Snell, K.~Swersky, and R.~Zemel.
\newblock Prototypical networks for few-shot learning.
\newblock In \emph{NeurIPS}, 2017.

\end{thebibliography}

\end{document}
